\section{Q1. Do benchmark gains transfer? Not reliably --- this panel specializes}
\label{sec:actI}

The first act answers the most basic question a practitioner has after fine-tuning a guard:
\emph{what did the fine-tune actually buy, and does it hold up off the training distribution?} The
leaderboard answer --- a single higher number --- silently blends two very different things: doing
better on benchmark \emph{sources the guard trained on} (``represented''), and doing better on
datasets it \emph{never saw} (``transfer''). Act I separates them with a strictly paired,
same-checkpoint measurement and finds a clean, repeatable pattern: supervised fine-tuning (SFT)
delivers a large, uniform represented-source gain and essentially no average transfer gain --- the
guard \emph{specializes}.

\begin{background}{What ``paired'' buys us, and why it matters}
Most guard papers compare \emph{different} models: model $X$'s guard against model $Y$'s guard. That
confounds ``the fine-tune helped'' with ``$X$ was a better starting model than $Y$.'' We instead hold
the checkpoint fixed and compare each guard \emph{against its own base}, on the identical rows, with
the identical single-token score head ($s(x)=z_{\text{unsafe}}-z_{\text{safe}}$; see
\Cref{sec:background}) and the identical tie-aware macro-AP. The reported quantity is therefore a
\emph{within-checkpoint change}, $\Delta = \text{SFT} - \text{base}$, not a cross-model gap. This is
the only way to attribute a movement to the fine-tune rather than to the luck of the starting point.
\end{background}

\subsection{The paired base-to-SFT design}
\label{sec:actI-design}

\paragraph{Panel and manifest.} All four checkpoints (Qwen2.5-1.5B, SmolLM2-1.7B, SmolLM3-3B,
Qwen3-4B) are fine-tuned from the identical frozen manifest of $1{,}200$ rows: $400$ from each of three
represented sources --- \code{toxicchat}, \code{prompt\_injections}, and
\code{jailbreak\_classification} --- balanced $200$ safe / $200$ unsafe within each source. The manifest
was decontaminated against the evaluation sets and its data order fixed by seed $42$, so every
checkpoint sees the same examples in the same order; the only thing that varies across the five runs
per checkpoint is the training seed.

\paragraph{Recipe.} Each guard is a LoRA adapter (\code{r=32}, \code{alpha=64}, \code{dropout=0.05})
attached to the attention and MLP projections (\code{q,k,v,o,gate,up,down}), trained for $300$ steps at
learning rate \code{2e-4} on a cosine schedule with $3\%$ warmup, effective batch size $4$, and maximum
sequence length $1024$. We run five seeds ($42$--$46$) per checkpoint, giving $20$ SFT guards in total.
Only the adapter is trained; the base weights are frozen, so ``base'' means the same model \emph{before}
the guard adapter, evaluated with the identical scoring head.

\paragraph{Two evaluation regimes.} Every guard and its base are scored in two disjoint regimes.
\emph{Represented} = held-back rows from the three training sources (same sources, unseen rows);
\emph{transfer} = four datasets whose rows were never used in training --- \code{jailbreakbench},
\code{xstest}, \code{wildguardtest}, and \code{wildjailbreak}. We additionally probe two single-class
stress sets that isolate the two failure directions: \code{OR-Bench} (benign prompts, to measure
over-refusal) and \code{HarmBench} (hard attack prompts, to measure missed catches).

\paragraph{Estimand and intervals.} We summarise each regime by macro-AP (average precision averaged
with equal weight across benchmarks, so a large benchmark cannot dominate) and put a two-sided $95\%$
interval on each change with a \emph{paired hierarchical bootstrap} that keeps every adapter tied to its
own base and respects the checkpoint$\to$seed grouping. The estimand is deliberately narrow: it is the
change \emph{conditional on this fixed, purposively-chosen four-checkpoint panel and this manifest},
which was inspected during development. These are estimation intervals, not significance tests, and the
result is \emph{retrospective, not confirmatory} --- no causal or population claim is licensed.

\subsection{The primary result: a uniform represented gain, a split transfer effect}
\label{sec:actI-primary}

\Cref{tab:primary} reports, per checkpoint, the untuned-base and mean-SFT macro-AP in each regime, the
paired base-to-SFT change with its interval, and the fixed-panel aggregate.

\begin{table}[htbp]\centering\small
\caption{Act I --- fixed-panel result: base and mean-SFT macro-AP by regime, paired base-to-SFT deltas
with two-sided $95\%$ hierarchical-bootstrap intervals, and the fixed-panel aggregate. Descriptive under
the locked precision-focused mode; conditional on this panel.}
\label{tab:primary}
\resizebox{\textwidth}{!}{% Auto-generated by analyze_paper_a_sft.py -- do not edit by hand.
\begin{tabular}{lrrrrrr}
\toprule
Checkpoint & Rep base & Rep SFT & $\Delta$ Rep [two-sided 95\% percentile CI] & Tr base & Tr SFT & $\Delta$ Tr [two-sided 95\% percentile CI] \\
\midrule
Qwen2.5-1.5B & 0.6334 & 0.9878 & 0.3544 [0.2731, 0.4150] & 0.8187 & 0.7798 & -0.0389 [-0.0829, 0.0062] \\
SmolLM2-1.7B & 0.4524 & 0.9806 & 0.5282 [0.4555, 0.5748] & 0.7904 & 0.8304 & 0.0400 [0.0003, 0.0776] \\
SmolLM3-3B & 0.6621 & 0.9751 & 0.3130 [0.2419, 0.3701] & 0.9102 & 0.8234 & -0.0869 [-0.1114, -0.0613] \\
Qwen3-4B & 0.8855 & 0.9837 & 0.0981 [0.0545, 0.1479] & 0.9438 & 0.7939 & -0.1499 [-0.1963, -0.1050] \\
\midrule
Fixed-panel aggregate & -- & -- & 0.3234 [0.2647, 0.3690] & -- & -- & -0.0589 [-0.0837, -0.0321] \\
\bottomrule
\end{tabular}
}
\end{table}

\paragraph{Represented: everyone reaches the same ceiling.} The represented-source gain is large and,
strikingly, \emph{uniform in its destination}: regardless of where the base started, every SFT guard
lands at macro-AP $\approx 0.98$. SmolLM2-1.7B climbs from a weak $0.4524$ to $0.9806$; Qwen2.5-1.5B
from $0.6334$ to $0.9878$; SmolLM3-3B from $0.6621$ to $0.9751$; and Qwen3-4B, already strong at
$0.8855$, to $0.9837$. Because the finish line is shared, the \emph{size} of the represented gain is
essentially dictated by how far below the ceiling the base sat --- a $+0.5282$ jump for SmolLM2 versus
only $+0.0981$ for Qwen3-4B. The fixed-panel aggregate is $\RepDelta$
$[\RepCILower, \RepCIUpper]$: SFT reliably teaches each model to rank the training sources' held-back
rows near-perfectly. Mechanistically this is unsurprising --- the manifest \emph{is} those sources, so
the adapter learns their label conventions and surface cues.

\paragraph{Transfer: a small average that hides opposite signs.} On data the guard never saw, the same
fine-tune does something quite different. The aggregate transfer change is only $\TransferDelta$
$[\TransferCILower, \TransferCIUpper]$ --- close to zero and slightly negative. But this panel average
is a mirage: it pools \emph{opposing} per-checkpoint effects. The change is positive for the weakest
base and strongly negative for the strongest, spanning
$+0.0400$ $[+0.0003, +0.0776]$ for SmolLM2-1.7B, $-0.0389$ $[-0.0829, +0.0062]$ for Qwen2.5-1.5B,
$-0.0869$ $[-0.1114, -0.0613]$ for SmolLM3-3B, and $-0.1499$ $[-0.1963, -0.1050]$ for Qwen3-4B. The
ordering tracks base transfer strength (\Cref{tab:primary}): the models with the most transfer skill to
begin with (SmolLM3 at $0.9102$, Qwen3-4B at $0.9438$) give the most of it back, while the model that
had the least ($0.7904$) is the only one that improves. In other words, SFT does not add a portable
notion of ``unsafe''; it reshapes the score toward the training sources, and for a model that already
generalised well that reshaping is a net loss off-distribution (\Cref{fig:act1-bars}).

\begin{figure}[t]\centering
\includegraphics[width=0.82\linewidth]{fig_act1_percheckpoint.pdf}
\caption{Per-checkpoint view of Act~I: SFT lifts represented-source macro-AP for every checkpoint (blue),
but the transfer change (orange) ranges from $+0.04$ (SmolLM2) to $-0.15$ (Qwen3-4B) --- the fixed-panel
average hides opposite signs.}
\label{fig:act1-bars}
\end{figure}

\begin{takeaway}{Why the $\TransferDelta$ transfer average is misleading}
It averages effects that point in opposite directions --- SmolLM2 $+0.0400$ against Qwen3-4B $-0.1499$.
The panel-level number is small; the panel itself is \emph{not} uniform, and reporting only the
aggregate would erase the most important finding (who loses, and how much). Read the per-checkpoint
column, not the bottom row.
\end{takeaway}

\subsection{Where the transfer losses live: a per-benchmark decomposition}
\label{sec:actI-perbench}

Aggregating over benchmarks can also hide \emph{which} data the guard forgets how to rank. The upper
block of \Cref{tab:sensitivity} decomposes the change benchmark by benchmark. On the represented side
every source rises sharply --- \code{toxicchat} $+0.1880$, \code{prompt\_injections} $+0.3704$, and
\code{jailbreak\_classification} $+0.4119$ --- confirming the gain is broad across the training sources,
not driven by one easy set.

On the transfer side the losses are \emph{concentrated on the jailbreak-style adversarial sets}:
\code{jailbreakbench} $-0.0776$, \code{wildjailbreak} $-0.0792$, and \code{wildguardtest} $-0.0669$ all
fall by a comparable amount, while \code{xstest} --- which contrasts genuinely unsafe prompts against
benign look-alikes rather than adversarial attacks --- barely moves at $-0.0120$. The pattern is
coherent with the mechanism above: the training sources over-represent a particular flavour of unsafe
text, so after tuning the guard ranks \emph{those} confidently but loses discrimination precisely on the
adversarial, distribution-shifted prompts that a deployed guard most needs to catch. Specialization is
not uniform forgetting; it is forgetting the hardest, most out-of-distribution cases first.

\subsection{The specialization plane: 15 of 20 guards specialize}
\label{sec:actI-plane}

The clearest single view of Act I is the \emph{specialization plane} (\Cref{fig:plane}), which plots
each of the $20$ (checkpoint, seed) guards by its represented change (horizontal) against its transfer
change (vertical). Four quadrants have direct meaning: lower-right = represented up, transfer down
(``specialize''); upper-right = both up (``uniform gain''); lower-left = both down (``uniform loss'');
upper-left = transfer-favoured. \SpecializationSeedCount{} of \TotalSeedCount{} guards land in the
specialize quadrant, \UniformGainSeedCount{} in uniform gain, and \emph{none} in uniform loss or the
transfer-favoured quadrant. So specialization is the dominant --- but not universal --- outcome: the
five uniform-gain points are four of the five SmolLM2 seeds \emph{plus} Qwen2.5's seed~42 --- concentrated
on the two checkpoints with the weakest base transfer, which have the most room to improve there (SmolLM2's
remaining seed, 46, narrowly specializes at $-0.001$). The absence of any uniform-loss point matters: SFT never simply degrades the
guard everywhere; it trades transfer for represented ranking. Per-seed values behind the plane are
tabulated in \Cref{tab:seed-values}, and show the effect is stable within each checkpoint (e.g. every
Qwen3-4B seed is a substantial transfer loss, every SmolLM2 seed a represented gain).

\begin{figure}[htbp]\centering
\includegraphics[width=0.6\textwidth]{specialization_plane.pdf}
\caption{Each point is one (checkpoint, seed): horizontal axis is the represented-source macro-AP
change, vertical axis the transfer change. \SpecializationSeedCount{}/\TotalSeedCount{} land in the
lower-right specialization quadrant --- the visual signature of Act I.}
\label{fig:plane}
\end{figure}

\begin{table}[htbp]\centering\scriptsize\setlength{\tabcolsep}{2.5pt}
\caption{Observed represented-source and transfer macro-AP change for every SFT seed (Act~I),
generated from the same validated score matrix as \Cref{tab:primary}. Every Qwen3-4B seed is a substantial
transfer loss and every SmolLM2-1.7B seed a represented gain --- the per-seed stability behind the
specialization plane (\Cref{fig:plane}).}
\label{tab:seed-values}
% Auto-generated by analyze_paper_a_sft.py -- do not edit by hand.
\begin{tabular}{lrrr}
\toprule
Checkpoint & Seed & $\Delta$ represented AP & $\Delta$ transfer AP \\
\midrule
Qwen2.5-1.5B & 42 & 0.3569 & 0.0082 \\
Qwen2.5-1.5B & 43 & 0.3535 & -0.0261 \\
Qwen2.5-1.5B & 44 & 0.3506 & -0.0545 \\
Qwen2.5-1.5B & 45 & 0.3558 & -0.0438 \\
Qwen2.5-1.5B & 46 & 0.3551 & -0.0786 \\
SmolLM2-1.7B & 42 & 0.5304 & 0.0805 \\
SmolLM2-1.7B & 43 & 0.5273 & 0.0720 \\
SmolLM2-1.7B & 44 & 0.5253 & 0.0250 \\
SmolLM2-1.7B & 45 & 0.5309 & 0.0238 \\
SmolLM2-1.7B & 46 & 0.5272 & -0.0013 \\
SmolLM3-3B & 42 & 0.3045 & -0.0812 \\
SmolLM3-3B & 43 & 0.3253 & -0.0925 \\
SmolLM3-3B & 44 & 0.3051 & -0.0629 \\
SmolLM3-3B & 45 & 0.3146 & -0.0923 \\
SmolLM3-3B & 46 & 0.3156 & -0.1055 \\
Qwen3-4B & 42 & 0.0953 & -0.1222 \\
Qwen3-4B & 43 & 0.0936 & -0.1536 \\
Qwen3-4B & 44 & 0.1038 & -0.0949 \\
Qwen3-4B & 45 & 0.0997 & -0.2291 \\
Qwen3-4B & 46 & 0.0981 & -0.1497 \\
\bottomrule
\end{tabular}

\end{table}

\subsection{At a deployable operating point}
\label{sec:actI-op}

Ranking is threshold-free, but deployment is not: at some point you must pick a cutoff and turn scores
into \code{block}/\code{allow} decisions. To see the deployment consequence of specialization we select
each guard's threshold to hit a $5\%$ false-positive-rate (FPR) target on a calibration split, then read
off the realized rates. The lower blocks of \Cref{tab:sensitivity} report them.

\paragraph{Represented recall soars.} On the represented sources, catching the unsafe prompts (TPR)
jumps from \RepBaseTPRPct\% (base) to \RepSFTTPRPct\% (SFT) --- the operating-point face of the same
$\approx 0.98$ ranking gain --- while the represented false-alarm rate stays low and even edges down
(\RepBaseFPRPct\% to \RepSFTFPRPct\%). If you only measured on the training sources, the guard would
look strictly and dramatically better.

\paragraph{Transfer pays for it.} Off-source the picture inverts. The transfer benchmark-macro FPR
climbs from \TransferBaseFPRPct\% (base) to \TransferSFTFPRPct\% (SFT) --- the tuned guard raises nearly
twice as many false alarms on data it never trained on --- and the pooled-negative FPR blows past target
from \TransferBasePooledFPRPct\% to \TransferSFTPooledFPRPct\%, so the calibrated threshold simply does
not transfer. Transfer recall rises only modestly (\TransferBaseTPRPct\% to \TransferSFTTPRPct\%),
nowhere near the represented jump --- and that rise is \emph{not} iso-FPR, which matters more than its
size. The two recalls are realized at \TransferBaseFPRPct\% and \TransferSFTFPRPct\% macro FPR, so the
tuned guard buys much of its extra recall by alarming roughly twice as often; a recall comparison at
unequal false-alarm rates is not a comparison of discriminative power.

\paragraph{The fair comparison: an equal false-alarm budget.} So we make the alarm budgets equal and
re-read the same rows. \Cref{tab:matchedfpr} gives each tuned guard the threshold at which its pooled
transfer false-alarm rate \emph{matches its own base's}, and then compares recalls. The modest apparent
gain does not merely shrink --- it reverses, on \emph{all four} checkpoints and by a wide margin: transfer recall goes \MatchedTransferBase{}\,$\to$\,\MatchedTransferSft{}
(\MatchedTransferDelta) and \texttt{HarmBench} recall \MatchedHarmBase{}\,$\to$\,\MatchedHarmSft{}
(\MatchedHarmDelta). At an equal alarm budget the tuned guard catches \emph{less than half} of what its
own untuned base catches off-source. The direction is stable across the three quantile conventions we
tried (panel mean \MatchedStabilityLo{} to \MatchedStabilityHi).

This is the honest form of the operating-point result, and it is the one the rest of Act~I predicts: the
$+0.06$ recall gain in \Cref{tab:sensitivity} is an artifact of comparing at unequal alarm rates. It needs
no GPU and no pinned environment --- matching false-alarm rates is ranking arithmetic on the same
committed \code{score\_raw} and \code{gold} columns every other covered artifact uses, so it is
regenerated and byte-checked by \code{make verify} (\Cref{sec:repro}); the emitter reproduces every
published row of \Cref{tab:sensitivity} exactly before adding the matched column, which is how we know the
two tables read the same underlying scores. Earlier drafts described this reconstruction as a
\emph{direction} that required the lock-pinned environment. That was wrong on both counts, and the
measured result is both larger and sharper than the hedge it replaces.

The single-class stress sets make the cost concrete: on the hard
\code{HarmBench} attacks the tuned guard's recall \emph{falls} from \HarmBenchBaseRecallPct\% to
\HarmBenchSFTRecallPct\% --- it catches less of the very content a guard exists to stop --- while
benign over-refusal on \code{OR-Bench} is essentially flat (\ORBenchBaseFPRPct\% to \ORBenchSFTFPRPct\%).
Since \code{OR-Bench} is \emph{also} unseen data, ``off-distribution'' cannot be what distinguishes it:
the extra false alarms are specific to the transfer suite's negatives, and we do not identify what
separates them. What the flat OR-Bench line does rule out is a blanket increase in caution.

Note that this \code{HarmBench} drop needs no threshold caveat at all: the tuned guard catches
\emph{less} while alarming \emph{more} (\TransferSFTFPRPct\% vs \TransferBaseFPRPct\% macro transfer
FPR), so it is dominated --- worse on both axes at once, not traded off. Equalising the budget only
widens the gap, to \MatchedHarmSft{} against the base's \MatchedHarmBase{} (\Cref{tab:matchedfpr}).

% tabcolsep was 6pt, which pushed this generated tabular 4.7pt past the text block.
\begin{table}[htbp]\centering\small\setlength{\tabcolsep}{4.5pt}
\caption{Act I --- per-benchmark paired deltas (upper) and the calibration-targeted $5\%$ FPR operating
point plus single-class stress diagnostics (lower). Generated from the lock-bound scores.}
\label{tab:sensitivity}
% Auto-generated by analyze_paper_a_sft.py -- do not edit by hand.
\begin{tabular}{llrrrr}
\toprule
Regime / benchmark & Metric & Base & SFT & $\Delta$ & $N$ \\
\midrule
represented / toxicchat & AP & -- & -- & 0.1880 & -- \\
represented / prompt\_injections & AP & -- & -- & 0.3704 & -- \\
represented / jailbreak\_classification & AP & -- & -- & 0.4119 & -- \\
transfer / jailbreakbench & AP & -- & -- & -0.0776 & -- \\
transfer / xstest & AP & -- & -- & -0.0120 & -- \\
transfer / wildguardtest & AP & -- & -- & -0.0669 & -- \\
transfer / wildjailbreak & AP & -- & -- & -0.0792 & -- \\
\midrule
represented & TPR@target FPR & 0.1296 & 0.7686 & 0.6390 & 313 \\
represented & benchmark-macro realized FPR & 0.0169 & 0.0108 & -0.0061 & 364 \\
represented & pooled-negative realized FPR & 0.0213 & 0.0194 & -0.0019 & 364 \\
transfer & TPR@target FPR & 0.5171 & 0.5807 & 0.0636 & 790 \\
transfer & benchmark-macro realized FPR & 0.0805 & 0.1546 & 0.0741 & 790 \\
transfer & pooled-negative realized FPR & 0.0430 & 0.1704 & 0.1273 & 790 \\
\midrule
Stress / OR-Bench & benign FPR & 0.1181 & 0.1201 & 0.0020 & 400 \\
Stress / HarmBench & recall & 0.7800 & 0.6003 & -0.1797 & 200 \\
\bottomrule
\end{tabular}

\end{table}

% GENERATED by reproduce.py (matched_fpr.py) from artifacts/paper_a_sft_v2/scores/scores.parquet
\begin{table}[htbp]\centering\small\setlength{\tabcolsep}{4.5pt}
\caption{\textbf{The same operating point read at an equal false-alarm budget.} \Cref{tab:sensitivity} compares recalls at each guard's \emph{own} calibrated threshold, where the tuned guard alarms far more often (17.0\% vs 4.3\% pooled transfer FPR) --- so it buys recall with alarms, and the two recalls are not comparable. Here each SFT seed is instead thresholded so its pooled transfer false-alarm rate \emph{matches its own base's} (the budget column), which makes the recalls directly comparable. At an equal budget the tuned guard is worse on \textbf{all four} checkpoints and on both instruments: transfer recall 0.517$\rightarrow$0.217 (-0.300) and \texttt{HarmBench} recall 0.780$\rightarrow$0.203 (-0.577). The direction is stable across the three quantile conventions we tried (panel-mean transfer delta -0.300 to -0.290). Same committed rows and same scorer as \Cref{tab:sensitivity}; only the threshold rule changes, so this needs no GPU and is byte-checked by \code{make verify}. \textbf{This is a retrospective ROC point, not a deployable threshold}: the quantile is read off the \emph{same} labelled negatives the recall is then measured on, so a production system without labels could not place it. Read the row as ``recall at an empirical matched-FPR ROC point'', and see \Cref{sec:matched-fpr-limits} for what an operational version would require.}
\label{tab:matchedfpr}
\begin{tabular}{lrrrrrrr}\toprule
& FPR & \multicolumn{3}{c}{transfer recall (macro)} & \multicolumn{3}{c}{\texttt{HarmBench} recall} \\
\cmidrule(lr){3-5}\cmidrule(lr){6-8}
Checkpoint & budget & base & SFT own thr. & SFT matched & base & SFT own thr. & SFT matched \\
\midrule
Qwen2.5-1.5B & 7.7\% & 0.485 & 0.571 & 0.296 & 0.810 & 0.606 & 0.312 \\
SmolLM2-1.7B & 2.4\% & 0.416 & 0.588 & 0.221 & 0.610 & 0.615 & 0.224 \\
SmolLM3-3B & 1.5\% & 0.512 & 0.596 & 0.141 & 0.740 & 0.611 & 0.116 \\
Qwen3-4B & 5.6\% & 0.655 & 0.569 & 0.211 & 0.960 & 0.569 & 0.160 \\
\midrule
Panel mean & 4.3\% & 0.517 & 0.581 & \textbf{0.217} & 0.780 & 0.600 & \textbf{0.203} \\
\bottomrule
\end{tabular}
\end{table}


\begin{takeaway}{The deployment cost behind the ranking story}
Read together, the operating point says the specialized guard \emph{catches more of what it trained on,
raises more false alarms on what it didn't, and misses more hard attacks}. A leaderboard number computed
on represented sources would advertise only the first of these three.

And the apparent consolation --- ``at least transfer recall went up'' --- does not survive a fair
comparison. That $+0.06$ was bought with alarms. Give the tuned guard the same alarm budget as its own
base and its transfer recall drops on all four checkpoints --- the panel mean roughly \emph{halves}
(\MatchedTransferBase{} to \MatchedTransferSft{}), and per checkpoint the fall runs from $39\%$ of the
base's recall (Qwen2.5-1.5B) to $72\%$ of it (SmolLM3-3B). The threshold that looked safe in calibration is not the threshold you get in the field ---
a theme Acts~II and~III return to.
\end{takeaway}

\noindent On this four-checkpoint panel, post-SFT scores cluster near a benchmark-fixed endpoint, so ``stronger bases specialize more'' is largely arithmetic rather than skill; the full attractor analysis and \Cref{fig:attractor} are in \Cref{app:detail}.

\subsection{The whole ranking is not the operating region}
\label{sec:actI-lowfpr}

\Cref{sec:actI-op} equalised the alarm \emph{budget} and re-read one operating point. A more basic
question sits underneath it, and the report has not answered it: our headline metric is macro-AP,
an average of precision over the \emph{entire} ranking, while every deployment sentence we write is
about a guard that fires on a few percent of traffic. Those are different quantities. AP rewards
ordering deep in the negative mass --- the region below any threshold an inline guard would ever
use --- so a change in AP need not imply the same change where the guard is actually placed.

\Cref{tab:lowfpr} settles it by recomputing the identical eight cells under three metrics on the
\emph{same} committed rows: macro-AP, the benchmark-macro one-way partial AUC over
FPR~$[0,\LowFprBudget]$, and benchmark-macro recall at that budget. Partial AUC integrates the ROC
only inside the alarm budget and is normalised to the mean TPR there, so it reads on the same scale
as a recall --- against a chance floor of $\LowFprChance$, not $0.5$. Intervals are the report's
usual paired bootstrap over \LowFprNBoot{} replicates, resampling evaluation \code{family\_id}
clusters and training seeds.

% GENERATED by reproduce.py (low_fpr.py) from artifacts/paper_a_sft_v2/scores/scores.parquet
\begin{table}[htbp]\centering\small\setlength{\tabcolsep}{5pt}
\caption{\textbf{The same eight cells, read in the operating region instead of over the whole ranking.} Paired base$\to$SFT change on identical committed rows under three metrics: benchmark-macro AP (the report's primary metric), benchmark-macro one-way partial AUC over FPR $[0,0.05]$ (mean TPR inside the alarm budget; chance floor $0.025$, not $0.5$), and benchmark-macro TPR at that budget. Brackets are two-sided $95\%$ paired bootstrap intervals over 2,000 replicates resampling evaluation \code{family\_id} clusters and training seeds, the same protocol as the rest of the report. \textbf{No sign flips}: every cell moves the same way under all three metrics, so Act~I's direction survives being read at a deployable operating point. What does not survive is the \emph{size}. macro-AP understates both halves of the trade --- the represented gain is +0.323 on AP against +0.686 on pAUC ($2.1\times$), and the transfer cost is -0.059 against -0.174 ($3.0\times$), with panel-mean budget recall moving -0.199. Averaging precision over the whole ranking credits ordering in the deep negative mass, where an inline guard never operates. Same rows, same scorer, no GPU: only the metric changes, and it is byte-checked by \code{make verify}.}
\label{tab:lowfpr}
\begin{tabular}{@{}lccc@{}}\toprule
Checkpoint & $\Delta$ macro-AP & $\Delta$ pAUC$[0,0.05]$ & $\Delta$ TPR@0.05 \\
\midrule
\multicolumn{4}{@{}l}{\emph{Represented (\code{id\_test})}} \\
\quad Qwen2.5-1.5B & +0.354~[+0.272,\,+0.412] & +0.802~[+0.723,\,+0.863] & +0.802~[+0.703,\,+0.860] \\
\quad SmolLM2-1.7B & +0.528~[+0.456,\,+0.572] & +0.791~[+0.731,\,+0.851] & +0.842~[+0.771,\,+0.889] \\
\quad SmolLM3-3B & +0.313~[+0.242,\,+0.367] & +0.676~[+0.567,\,+0.767] & +0.712~[+0.550,\,+0.782] \\
\quad Qwen3-4B & +0.098~[+0.055,\,+0.148] & +0.475~[+0.305,\,+0.563] & +0.357~[+0.195,\,+0.514] \\
\quad \textbf{Panel mean} & \textbf{+0.323} & \textbf{+0.686} & \textbf{+0.678} \\
\midrule
\multicolumn{4}{@{}l}{\emph{Transfer (\code{transfer\_test})}} \\
\quad Qwen2.5-1.5B & -0.039~[-0.082,\,+0.007] & -0.002~[-0.126,\,+0.085] & -0.081~[-0.198,\,+0.075] \\
\quad SmolLM2-1.7B & +0.040~[+0.001,\,+0.077] & +0.011~[-0.093,\,+0.118] & -0.012~[-0.108,\,+0.113] \\
\quad SmolLM3-3B & -0.087~[-0.111,\,-0.062] & -0.295~[-0.363,\,-0.219] & -0.265~[-0.364,\,-0.196] \\
\quad Qwen3-4B & -0.150~[-0.197,\,-0.106] & -0.409~[-0.499,\,-0.312] & -0.437~[-0.528,\,-0.311] \\
\quad \textbf{Panel mean} & \textbf{-0.059} & \textbf{-0.174} & \textbf{-0.199} \\
\bottomrule
\end{tabular}
\end{table}


\paragraph{The direction survives; the magnitude does not.} \LowFprSignFlips{} of the eight cells
change sign: every checkpoint moves the same way under all three metrics, so Act~I's qualitative
finding is not an artifact of choosing an average-ranking metric. What macro-AP does do is
systematically \emph{understate} the trade, on both sides. The represented gain is
$\LowFprRepApDelta$ on AP against $\LowFprRepPAucDelta$ on pAUC
($\LowFprRepAmplification\times$), and the transfer cost is $\LowFprTransApDelta$ against
$\LowFprTransPAucDelta$ ($\LowFprTransAmplification\times$), with panel-mean budget recall moving
$\LowFprTransTprDelta$. The worst cell is the recurring one: \LowFprWorstName{} gives back
$\LowFprWorstAp$ of macro-AP and $\LowFprWorstPAuc$ \LowFprWorstPAucCI{} of pAUC ---
$\LowFprWorstTpr$ in budget recall, which is to say roughly two of every five unsafe prompts it
would have caught at a $\LowFprBudget$ alarm rate before tuning.

One qualification cuts the other way, and it weakens a claim we make elsewhere. SmolLM2-1.7B is the
single checkpoint whose transfer \emph{improves} under SFT on macro-AP ($+0.0400$, an interval that
clears zero). In the operating region that gain does not survive: pAUC $+0.011$ and budget recall
$-0.012$, both straddling zero. So ``positive for the weakest base'' is a statement about average
ranking, not about deployable recall, and \Cref{sec:actI-primary}'s ``the model that had the least
is the only one that improves'' should be read with that scope.

\begin{takeaway}{Why we report both, and which one to act on}
macro-AP is the right instrument for the question Act~I asks --- \emph{did the fine-tune reorder the
data?} --- and it is what makes the paired estimand comparable across benchmarks of different base
rates. It is the wrong instrument for the question a practitioner asks: \emph{what happens at the
alarm rate I can afford?} On this panel the two agree in sign and disagree by a factor of two to
three in size, always in the direction that makes specialization look milder than it is. Read the
AP columns to compare methods; read the pAUC and budget-recall columns before believing a
deployment number.
\end{takeaway}

\noindent This is a re-reading of evidence the report already has, at the same evidence tier: same
fixed panel, same inspected sources, same retrospective status, no new data and no GPU. It licenses
nothing that \Cref{tab:primary} did not already license --- it prices it differently.

\subsection{A recipe control: does anti-forgetting (KL-regularized) SFT preserve transfer?}
\label{sec:actI-klsft}

The specialization above was measured under \emph{one} fine-tuning recipe --- completion-only LoRA-SFT with
no explicit anchor to the base checkpoint. A fair objection is that the transfer loss might be a property of
this \emph{unregularized} recipe rather than of fine-tuning as such. The standard remedy is an
\emph{anti-forgetting} penalty that keeps the fine-tune close to the base in output space, so we add exactly
that --- a KL term on the SFT loss:
\begin{equation}
\mathcal{L} \;=\; \underbrace{\mathrm{CE}(\text{verdict})}_{\text{Act~I recipe}} \;+\;
\beta\,\mathrm{KL}\!\big(\pi_\theta(\cdot\mid x)\,\big\|\,\pi_{\text{base}}(\cdot\mid x)\big),
\label{eq:klsft}
\end{equation}
evaluated on the completion tokens, with the frozen base $\pi_{\text{base}}$ recovered from the adapter's own
disabled path (no second model in memory, no base retraining). Everything else is held at the Act~I settings
--- same manifest, seeds, panel, scorer, and represented/transfer splits --- and $\beta=0$ reproduces vanilla
SFT \emph{exactly} (the cross-entropy is unchanged), so this is a strict one-knob generalization of the Act~I
recipe rather than a different method. We sweep $\beta\in\{\KLBetas\}$ over all four checkpoints and
\KLNSeeds\ seeds, and score every adapter through the identical single-token margin used everywhere else.

\input{generated/tab_klsft_gen}

\Cref{tab:klsft} reports transfer and represented macro-AP for the base, vanilla SFT ($\beta=0$, trained in
the \emph{same} environment as the KL runs so the contrast carries no hardware confound), and
KL-regularized SFT at each $\beta$. The anti-forgetting question is read from the transfer block: relative to
vanilla SFT, KL at $\beta=0.5$ changes transfer macro-AP by \KLBetaHalfTransferDelta\ on average across the
four checkpoints (at a represented cost of \KLBetaHalfRepDelta), and at $\beta=1.0$ by
\KLBetaOneTransferDelta\ (represented cost \KLBetaOneRepDelta). \KLVerdictSentence

\paragraph{An accidental noise floor, and why it bounds several claims in this report.} That
same-environment $\beta=0$ arm is also the closest thing we have to a \emph{repeat} of Act~I: identical
recipe, identical manifest, identical seeds, identical scorer --- only the execution environment differs.
It does not land on the same number. Comparing the SFT transfer column of \Cref{tab:klsft} with
\Cref{tab:primary} gives gaps of $0.014$, $0.009$, $0.009$ and $0.029$ (Qwen2.5, SmolLM2, SmolLM3,
Qwen3-4B), i.e.\ a mean of $0.015$ and a worst case of $0.029$ macro-AP for a recipe we intended to be
deterministic.

We report this because it is a ceiling on what any small effect in this report can mean. Effects
comfortably above it --- Act~I's represented gain ($+0.32$), the matched-budget recall collapse
($\MatchedTransferDelta$), composition's recovery over SFT ($+0.076$) --- are not threatened by it.
Effects at or below it should be read as unresolved: composition's $+0.017$ aggregate edge over the
\emph{base}, and the $\beta=1.0$ transfer change for SmolLM2 ($+0.004$), are both inside this envelope.
The bootstrap intervals elsewhere in the report resample evaluation rows and seeds; they do \emph{not}
capture this environment term, so they are narrower than a full reproduction would be.

\paragraph{The dial, priced in the operating region.} \Cref{sec:actI-lowfpr} showed macro-AP
understates Act~I's trade; it understates this one too, and asymmetrically.
\Cref{tab:lowfpr-kl} recomputes the same KL cells over FPR~$[0,\LowFprBudget]$. The direction is
unchanged at both $\beta$ --- KL buys transfer and charges represented ranking --- but at
$\beta{=}\LowFprKlBeta$ the transfer gain is $\LowFprKlTransApDelta$ on AP against
$\LowFprKlTransPAucDelta$ on pAUC ($\LowFprKlTransAmplification\times$) and
$\LowFprKlTransTprDelta$ in budget recall, while the represented cost is $\LowFprKlRepApDelta$
against $\LowFprKlRepPAucDelta$ ($\LowFprKlRepAmplification\times$). Both halves grow, and the
\emph{cost} half grows more than twice as fast as the benefit half. A dial that looks like
``two points of represented AP for six points of transfer AP'' is, where a guard is placed, closer
to twenty-one points of represented pAUC for fifteen. That does not make KL useless --- it makes the
choice sharper, and more clearly a choice.

% GENERATED by reproduce.py (low_fpr.py) from artifacts/klsft_v1/scores/
\begin{table}[htbp]\centering\small\setlength{\tabcolsep}{5pt}
\caption{\textbf{The KL dial, priced in the operating region.} Change from the in-environment $\beta{=}0$ arm, on identical committed rows, under the same three metrics as \Cref{tab:lowfpr}. The direction of \Cref{tab:klsft} is unchanged --- KL buys transfer and charges represented ranking at both $\beta$ --- but the exchange rate is not what macro-AP shows. At $\beta{=}0.5$ the transfer gain is +0.061 on AP against +0.149 on pAUC ($2.4\times$) and +0.163 in budget recall, while the represented cost is -0.035 on AP against -0.214 on pAUC ($6.2\times$). Both halves of the dial are larger where a guard is placed, and the represented half grows faster than the transfer half. This also sharpens \Cref{sec:adaptation}'s registered failure: read on AP the represented cost misses the $-0.02$ non-inferiority margin by under a factor of two, but read at the budget it misses by roughly 11$\times$. Point estimates only --- these are seed means over the same four checkpoints, not a re-run of the registered study's family bootstrap.}
\label{tab:lowfpr-kl}
\begin{tabular}{@{}lccc@{}}\toprule
Checkpoint & $\Delta$ macro-AP & $\Delta$ pAUC$[0,0.05]$ & $\Delta$ TPR@0.05 \\
\midrule
\multicolumn{4}{@{}l}{\emph{Transfer (\code{transfer\_test})}} \\
\quad Qwen2.5-1.5B ($\beta{=}0.5$) & +0.056 & +0.071 & +0.130 \\
\quad SmolLM2-1.7B ($\beta{=}0.5$) & +0.020 & +0.093 & +0.076 \\
\quad SmolLM3-3B ($\beta{=}0.5$) & +0.088 & +0.263 & +0.246 \\
\quad Qwen3-4B ($\beta{=}0.5$) & +0.081 & +0.170 & +0.200 \\
\quad \textbf{Panel mean} ($\beta{=}0.5$) & \textbf{+0.061} & \textbf{+0.149} & \textbf{+0.163} \\
\quad Qwen2.5-1.5B ($\beta{=}1$) & +0.045 & +0.062 & +0.119 \\
\quad SmolLM2-1.7B ($\beta{=}1$) & +0.005 & +0.070 & +0.080 \\
\quad SmolLM3-3B ($\beta{=}1$) & +0.085 & +0.258 & +0.237 \\
\quad Qwen3-4B ($\beta{=}1$) & +0.071 & +0.137 & +0.173 \\
\quad \textbf{Panel mean} ($\beta{=}1$) & \textbf{+0.051} & \textbf{+0.132} & \textbf{+0.152} \\
\midrule
\multicolumn{4}{@{}l}{\emph{Represented (\code{id\_test})}} \\
\quad Qwen2.5-1.5B ($\beta{=}0.5$) & -0.008 & -0.027 & -0.029 \\
\quad SmolLM2-1.7B ($\beta{=}0.5$) & -0.052 & -0.317 & -0.201 \\
\quad SmolLM3-3B ($\beta{=}0.5$) & -0.057 & -0.385 & -0.189 \\
\quad Qwen3-4B ($\beta{=}0.5$) & -0.022 & -0.126 & -0.073 \\
\quad \textbf{Panel mean} ($\beta{=}0.5$) & \textbf{-0.035} & \textbf{-0.214} & \textbf{-0.123} \\
\quad Qwen2.5-1.5B ($\beta{=}1$) & -0.010 & -0.046 & -0.030 \\
\quad SmolLM2-1.7B ($\beta{=}1$) & -0.072 & -0.380 & -0.265 \\
\quad SmolLM3-3B ($\beta{=}1$) & -0.091 & -0.505 & -0.364 \\
\quad Qwen3-4B ($\beta{=}1$) & -0.041 & -0.204 & -0.133 \\
\quad \textbf{Panel mean} ($\beta{=}1$) & \textbf{-0.053} & \textbf{-0.284} & \textbf{-0.198} \\
\bottomrule
\end{tabular}
\end{table}


\noindent It also sharpens the one preregistered verdict in this report. \Cref{sec:adaptation}'s RQ2
fails a registered $-0.02$ non-inferiority margin on represented macro-AP. Read on AP that margin is
missed by about $\LowFprKlMarginMultipleAp\times$; read in the operating region on this panel it is
missed by about $\LowFprKlMarginMultiplePAuc\times$. The registered criterion was not narrowly
missed. (These are point estimates on the four general checkpoints, not a re-run of the registered
study's own family bootstrap on its ten-checkpoint panel, so they qualify that verdict's magnitude
rather than restating its interval.)

\begin{takeaway}{What the recipe control shows}
\KLTakeaway
\end{takeaway}

\subsection{The deployment base rate re-spaces the ranking}
\label{sec:actI-prevalence}

One more axis quietly co-produces the score, and it is not the model at all: the \emph{prevalence} of
unsafe prompts. All the AP numbers above are measured on balanced (or near-balanced) pools, but real
inbound traffic is overwhelmingly benign --- unsafe prompts might be $1\%$ of requests. Average precision
depends on that base rate, and the dependence is exact: given a guard's ranking (its ROC), AP at any
prevalence $\pi_+$ is a fixed recomputation,
\begin{equation}
\mathrm{AP}(\pi_+) \;=\; \int_0^1 \frac{\pi_+\,u}{\pi_+\,u + (1-\pi_+)\,\mathrm{FPR}(u)}\,\mathrm{d}u,
\label{eq:ap-prevalence}
\end{equation}
where $u$ is recall and $\mathrm{FPR}(u)$ is the guard's false-positive rate at that recall --- so no new
model runs are needed, only the committed per-row scores.

\begin{figure}[htbp]\centering
\includegraphics[width=0.82\linewidth]{fig_prevalence.pdf}
\caption{\textbf{The prevalence re-spaces the ranking} (measured: \Cref{eq:ap-prevalence} applied to the
four base guards' committed transfer ROC, macro-averaged over the transfer sources). As unsafe prompts get
rarer, every guard's AP falls --- the strongest ranker (Qwen3-4B) drops from $0.94$ at balance to
$\approx0.56$ at $1\%$ prevalence, and the guard that ends lowest (Qwen2.5-1.5B) from $0.82$ to $0.11$ --- \emph{and} the
lower half re-orders (see text).}
\label{fig:prevalence}
\end{figure}

Two consequences follow, both recomputed exactly from the committed scores. First, the balanced AP is an
\emph{optimistic} reading of deployed precision: SmolLM3-3B's base transfer ranking scores $\mathrm{AP}=0.91$
on a balanced pool but only $\approx0.51$ at $1\%$ prevalence (\Cref{fig:prevalence}). Second --- and this is
the same lesson at the metric's own prior, though it is a re-spacing rather than a benchmark gain that
fails to transfer --- \textbf{low prevalence collapses and re-spaces the ranking}. The extremes are stable (Qwen3-4B stays first, SmolLM2 and Qwen2.5 stay in the lower half at every
prevalence), but the \emph{lower two re-order}: Qwen2.5-1.5B edges SmolLM2-1.7B at balance ($0.82$ vs.\ $0.79$)
yet falls behind it once positives are rare ($0.11$ vs.\ $0.21$ at $1\%$), because a scarce positive class
re-expands exactly the low-recall precision differences the balanced pool compresses. This is a re-spacing and
a partial re-order, not a wholesale winner-flip --- but it is enough that a single balanced AP can invert two
guards' apparent order at deployment prevalence. Reporting each headline guard as a curve $\mathrm{AP}(\pi_+)$
rather than a single balanced number is a zero-cost recompute (roadmap), and it is the honest way to state a
deployment precision.

\edb{Represented AP $\RepDelta$ (LCB $\RepLCB$) vs.\ transfer $\TransferDelta$ (UCB $\TransferUCB$);
specialization in $\SpecializationSeedCount/\TotalSeedCount$ seeds (\Cref{tab:primary}). Read in the
deployment region the direction holds and the size roughly triples: transfer
$\LowFprTransPAucDelta$ pAUC$[0,\LowFprBudget]$, budget recall $\LowFprTransTprDelta$
(\Cref{tab:lowfpr}). A base-anchored
KL penalty ($\beta{=}0.5$) buys back transfer (\KLBetaHalfTransferDelta{} vs.\ SFT) at a represented cost
\KLBetaHalfRepDelta{} (\Cref{sec:actI-klsft}).}
{Never replace a base guard on represented AP alone; compare each tune to its \emph{own} base on
represented \emph{and} held-out sets, and if you must SFT while caring about OOD, consider the KL penalty
--- but as a dial, not a default: the preregistered test in \Cref{sec:adaptation} finds its
represented-source cost fails the non-inferiority margin, and it leaves the two hardest specializers below
their own base on transfer.}
{Retrospective on a fixed four-checkpoint panel with an inspected manifest (the lexical overlap audit is
run and clean; the embedding-space check is not --- \Cref{sec:decontamination}); a
description of this recipe on these runs, not a population or causal law.}
